\section{Introduction}
\label{sec:intro}

\subsection{Setting}

Data-processing pipelines rarely stay inside one engine. A single analyst workflow may read
Parquet through PyArrow, reshape it with pandas or Polars, push an aggregation into DuckDB or
SQLite, and execute a larger relational plan in DataFusion or chDB. These systems are commonly
treated as interchangeable executors of the same table semantics, and in the common case they
behave as such. At the boundaries, however, they disagree: null and NaN propagation,
negative-zero comparisons, ordering and top-$k$ tie-breaking, group-by with null keys, join
semantics under duplicate keys, integer/float casting, string collation, and Arrow layout
details such as slice offsets and chunking. When such a disagreement occurs, both engines
return a plausible result and no error is raised. The pipeline continues with silently wrong
data.

\subsection{Problem}

Detecting these \emph{silent wrong-result} bugs in this ecosystem is hard for four reasons.
First, the testing surface is heterogeneous: one workflow must be lowered into DataFrame API
calls, Arrow compute kernels, SQL, and query-engine plans, each with its own type system and
capability set. Second, directly comparing raw outputs is dominated by false positives, because
legitimate differences in dtype, index, ordering, and presentation are ubiquitous; a canonical
representation and a classification layer are prerequisites, not refinements. Third,
techniques developed for SQL DBMSs assume a single language and a single query planner and do
not transfer to lazy/eager DataFrame execution or Arrow in-memory layout. Fourth, a fuzzing
finding is not a bug: it becomes a contribution only after it is re-checked, minimized,
deduplicated by root cause, and independently confirmed upstream.

\subsection{Gap}

Prior work has advanced three strands separately. DBMS logic-bug testing provides powerful
oracle families---pivoted query synthesis, non-optimizing reference engines, query
partitioning, cardinality and metamorphic oracles---but targets SQL engines. DataFrame and
workload-generation work targets Python data-processing libraries, but emphasizes workload
realism and replay rather than a cross-model semantic oracle. Specialized data-system testing
targets graph, spatial, or time-series engines. To our knowledge no prior published system
unifies DataFrame APIs, Arrow processing, embedded SQL, and analytical query engines under one
typed testing surface with one normalizer, one compositional oracle, and one evidence
pipeline; we deliberately do not claim priority for differential, metamorphic, or DataFrame
testing as techniques.

\subsection{Thesis}

A single typed workflow IR with capability-aware lowering, a shared semantic normalizer, a
compositional multi-endpoint oracle complex, and an end-to-end evidence pipeline can find,
reproduce, and substantiate latest-version silent wrong-result bugs across the whole tabular
data-processing ecosystem, and can do so with controlled noise and auditable provenance.

\subsection{Contributions}

\begin{itemize}[leftmargin=*,itemsep=2pt]
  \item \textbf{C1 (unified testing surface).} A typed workflow IR and capability-aware
  lowering that place Python DataFrame APIs, Arrow compute, embedded SQL, and query engines on
  one testing surface (\S\ref{sec:approach-surface}), evaluated for reuse and per-family
  coverage in RQ6.
  \item \textbf{C2 (normalizer + compositional oracle).} A semantic normalizer and a
  multi-endpoint oracle complex that fuses differential checking, metamorphic relations, and
  deterministic latest-version probes, with derivation, applicability, and observation
  certificates and separate \textsc{satisfied}/\textsc{violated}/\textsc{inapplicable}/
  \textsc{inconclusive} verdicts (\S\ref{sec:approach-oracle}); evaluated for noise control
  (RQ2) and oracle complementarity (RQ3).
  \item \textbf{C6 (evidence pipeline).} An end-to-end pipeline from candidate to
  issue-ready root to independently confirmed family, with fresh/replay physical separation
  and a strict counting contract (\S\ref{sec:approach-evidence}); evaluated in RQ5.
  \item \textbf{Oracle core (C7).} A finite compositional semantic HyperContract compiled by
  forward-property inference and backward-observability analysis, driving a coverage-debt
  scheduler over a declared target lattice, with authority-bound replayable evidence
  (\S\ref{sec:approach-hypercontract}, \S\ref{sec:impl-authority}); evaluated in RQ-C7.
\end{itemize}

\subsection{Results preview}

Using \sys{} across DataFusion, Polars, DuckDB, and PyArrow, we report \textbf{nine strict
confirmed latest-version bug families} (seven upstream-labeled, two fixed upstream), with one
further finding pending independent confirmation and one matching a pre-existing upstream
issue reported separately (Table~\ref{tab:bugs}). \todo{Add regenerated RQ2/RQ3/RQ5/RQ6
summary sentences once W3 completes.}

\subsection{Artifact availability}

\todo{Describe the artifact package: frozen source commit, exact environment lock, seed
policy, raw journals, receipts, and the single reproduction script.}
